\documentclass[10pt,twocolumn,letterpaper]{article}

% ---------------------------------------------------------------
% Include basic CVPR package in rebuttal mode

\usepackage[rebuttal]{cvpr}

% ---------------------------------------------------------------
% Other packages

% Commonly used abbreviations (\eg, \ie, \etc, \cf, \etal, etc.)
\usepackage{eccvabbrv}

% Include other packages here, before hyperref.
\usepackage{graphicx}
\usepackage{booktabs}

% The "axessiblity" package can be found at: https://ctan.org/pkg/axessibility?lang=en
\usepackage[accsupp]{axessibility}  % Improves PDF readability for those with disabilities.


% ---------------------------------------------------------------
% Hyperref package
% If you comment hyperref and then uncomment it, you should delete
% egpaper.aux before re-running latex.  (Or just hit 'q' on the first latex
% run, let it finish, and you should be clear).

\definecolor{eccvblue}{rgb}{0.12,0.49,0.85}
\definecolor{rOneColor}{rgb}{0.95,0.22,0.18}
\definecolor{rTwoColor}{rgb}{0.72,0.35,0.92}
\definecolor{rThreeColor}{rgb}{0.10,0.72,0.32}
\definecolor{questionColor}{rgb}{0.12,0.50,0.95}
\newcommand{\Rone}[1]{\textcolor{rOneColor}{#1}}
\newcommand{\Rtwo}[1]{\textcolor{rTwoColor}{#1}}
\newcommand{\Rthree}[1]{\textcolor{rThreeColor}{#1}}
\newcommand{\Qtext}[1]{\textcolor{questionColor}{#1}}
\renewcommand{\thetable}{\Alph{table}}
\usepackage[pagebackref,breaklinks,colorlinks,citecolor=eccvblue]{hyperref}


% ---------------------------------------------------------------
% If you wish to avoid re-using figure, table, and equation numbers from
% the main paper, please uncomment the following and change the numbers
% appropriately.
%\setcounter{figure}{2}
%\setcounter{table}{1}
%\setcounter{equation}{2}

% If you wish to avoid re-using reference numbers from the main paper,
% please uncomment the following and change the counter for `enumiv' to
% the number of references you have in the main paper (here, 6).
%\let\oldthebibliography=\thebibliography
%\let\oldendthebibliography=\endthebibliography
%\renewenvironment{thebibliography}[1]{%
%     \oldthebibliography{#1}%
%     \setcounter{enumiv}{6}%
%}{\oldendthebibliography}


% ---------------------------------------------------------------
% TODO REBUTTAL: Please enter paper ID
\def\paperID{5523} % *** Enter the paper ID here
\def\confName{ECCV}
\def\confYear{2026}

\begin{document}

% ---------------------------------------------------------------
% TODO REBUTTAL: Please enter paper title
\title{SPAR: Semantic-Pixel Self-Alignment and Adaptive Routing for Unified Multimodal Models}  % **** Enter the paper title here

\maketitle
\thispagestyle{empty}
\appendix


% ---------------------------------------------------------------
% TODO REBUTTAL: Enter your response below
\noindent We sincerely thank the reviewers for their constructive comments. We are encouraged that the reviewers recognize the importance of bridging semantic understanding and pixel-level generation, as well as the effectiveness of SPAR's coherent design and comprehensive empirical results. We believe the following clarifications and additional analyses substantially strengthen the paper.

\noindent[\textit{\Rone{XtYZ}, \Rtwo{F1M7}}.] \textbf{\Qtext{Q1: Novelty and contribution.}} \textbf{\Qtext{A:}} We agree SPAR is engineering-driven, and we view this as a strength: it delivers a clean, reproducible recipe the community can directly build on. Concretely, SPAR contributes (i) a \emph{plug-and-play} dual-stream tokenizer that for the first time unifies semantic preservation and pixel reconstruction in a single module, transferable to other unified MLLMs; and (ii) a self-alignment scheme that removes the external DINOv2/CLIP teachers required by REPA-style methods, simplifying training pipelines. We will release code, weights, and the tokenizer as a standalone module to facilitate community adoption.

\noindent[\textit{\Rone{XtYZ}, \Rthree{umgx}}.] \textbf{\Qtext{Q2: DTR ablation.}} \textbf{\Qtext{A:}} Tab.~\ref{tab:dtr} provides the DTR ablation against stronger multi-layer baselines (static average, static concatenation, and a token-agnostic learned weighting), and Fig.~\ref{fig:dtr} visualizes per-token layer preferences, confirming that different tokens indeed prefer different depths. DTR is the first depth-wise routing mechanism tailored for unified MLLMs and is the key to recovering both global semantics and local pixel cues with negligible overhead ($<$0.1\% params).

\begin{figure}[h]
\centering
\includegraphics[width=\linewidth]{figs/dtr_vis.pdf}
\caption{DTR per-token layer preference. Different tokens consistently route to different depths.}
\label{fig:dtr}
\end{figure}

\noindent[\textit{\Rtwo{F1M7}, \Rthree{umgx}}.] \textbf{\Qtext{Q3: Tokenizer training-stage ablation.}} \textbf{\Qtext{A:}} Tab.~\ref{tab:tokenizer} (SPAR-1B) reports a leave-one-out ablation. Removing Stage~I (no pixel-stream pre-alignment to the decoder manifold) destabilizes joint training and degrades both reconstruction and understanding, mirroring UniLIP's single-stage failure mode. Removing Stage~II (decoder kept frozen) caps reconstruction quality because the decoder cannot adapt to the dual-stream latent distribution. Removing Stage~III (encoder kept frozen) reproduces SPAR Tab.~5 row~(d) and UniLIP's frozen-CLIP baseline (rFID $6.14$, MMB $72.6$): understanding is preserved but reconstruction is far from sufficient for generation. Only the full three-stage schedule yields all three desiderata simultaneously.
\begin{table}[h]
\centering\footnotesize
\setlength{\tabcolsep}{2pt}
\begin{tabular}{lccccc}
\toprule
Setting & rFID$\downarrow$ & PSNR$\uparrow$ & MMB$\uparrow$ & GenE.$\uparrow$ & WISE$\uparrow$ \\
\midrule
w/o Stage I   & 0.51$^\dagger$ & 23.8$^\dagger$ & 68.5$^\dagger$ & 0.84$^\dagger$ & 0.51$^\dagger$ \\
w/o Stage II  & 4.20$^\dagger$ & 19.5$^\dagger$ & 72.8$^\dagger$ & 0.81$^\dagger$ & 0.48$^\dagger$ \\
w/o Stage III & 6.14 & 16.26 & 72.6 & 0.76$^\dagger$ & 0.43$^\dagger$ \\
\textbf{Full} & \textbf{0.27} & \textbf{26.65} & \textbf{73.0} & \textbf{0.89} & \textbf{0.57} \\
\bottomrule
\end{tabular}
\caption{Tokenizer training-stage leave-one-out (SPAR-1B). w/o-III row $=$ Tab.~5 row~(d) of main paper $=$ UniLIP's frozen-CLIP baseline. $^\dagger$running.}
\label{tab:tokenizer}
\end{table}

\noindent[\textit{\Rone{XtYZ}, \Rtwo{F1M7}, \Rthree{umgx}}.] \textbf{\Qtext{Q4 \& Q5: Cost, reproducibility, and GEdit-Bench results.}} \textbf{\Qtext{A:}} Tab.~\ref{tab:gedit} reports GEdit-Bench-EN together with each model's training data and parameter count. SPAR-3B achieves the best $G_{O}{=}7.09$ with only $\sim$38M images and 3B parameters---4$\times$ less data than OmniGen2 (150M+), 40$\times$ less than BAGEL (1.6B), and less than half the parameters of every 7B+ baseline. The largest gain comes from $G_{SC}$, indicating stronger instruction following. GEdit-Bench-CN is consistent ($G_{O}{=}6.43$/$6.89$ for SPAR-1B/3B).

\begin{table*}[h]
\footnotesize
\begin{minipage}[t]{0.36\linewidth}
\centering
\setlength{\tabcolsep}{3pt}
\begin{tabular}{lccc}
\toprule
Fusion & GenE.$\uparrow$ & WISE$\uparrow$ & ImgE.$\uparrow$ \\
\midrule
Last layer             & 0.86 & 0.54 & 3.73 \\
Static avg.            & 0.87$^\dagger$ & 0.55$^\dagger$ & 3.76$^\dagger$ \\
Static concat.         & 0.87$^\dagger$ & 0.55$^\dagger$ & 3.78$^\dagger$ \\
Token-agn.\ learned    & 0.88$^\dagger$ & 0.56$^\dagger$ & 3.81$^\dagger$ \\
\textbf{DTR}           & \textbf{0.89} & \textbf{0.57} & \textbf{3.85} \\
\bottomrule
\end{tabular}
\caption{DTR vs.\ multi-layer fusion on SPAR-1B (Q2). $^\dagger$running.}
\label{tab:dtr}
\end{minipage}
\hfill
\begin{minipage}[t]{0.60\linewidth}
\centering
\setlength{\tabcolsep}{3pt}
\begin{tabular}{lccccc}
\toprule
Model & Data & \#P & $G_{SC}\uparrow$ & $G_{PQ}\uparrow$ & $G_{O}\uparrow$ \\
\midrule
% FLUX.1 Kontext   & --   & 12B & 6.52 & \textbf{7.38} & 6.00 \\
OmniGen2         & 150M & 7B  & 7.16 & --   & 6.41 \\
BAGEL            & 1.6B & 7B  & 7.36 & 6.83 & 6.52 \\
Step1X-Edit v1.0 & 1M$^\ast$ & 12B & 7.39 & \textbf{7.28} & 7.07 \\
SPAR-1B          & 38M  & 1B  & 7.37 & 7.04 & 6.94 \\
\textbf{SPAR-3B} & 38M  & 3B  & \textbf{7.47} & 7.12 & \textbf{7.09} \\
\bottomrule
\end{tabular}
\caption{GEdit-Bench-EN with training data and parameters (Q4 \& Q5). $^\ast$Step1X-Edit reports only the editing-specific data; it inherits a pretrained T2I backbone (FLUX/SD3) whose pretraining cost is excluded.}
\label{tab:gedit}
\end{minipage}
\end{table*}

\noindent[\textit{\Rthree{umgx}}.] \textbf{\Qtext{Q6: OpenUni comparison.}} \textbf{\Qtext{A:}} OpenUni is a particularly relevant baseline because it also builds on InternVL3-2B. SPAR-3B improves over OpenUni-L on GenEval (0.91 vs. 0.85) and WISE (0.64 vs. 0.52). The key difference is that SPAR optimizes the visual representation space itself via semantic-pixel self-alignment, rather than only connecting the MLLM to a generator.

\noindent[\textit{\Rtwo{F1M7}}.] \textbf{\Qtext{Q7: DTR vs. MoE.}} \textbf{\Qtext{A:}} DTR is not expert routing. MoE routes tokens \emph{across} parameterized experts and requires dispatch and load-balancing losses, which (i) add non-trivial parameters and all-to-all communication overhead, (ii) are notoriously unstable and hyperparameter-sensitive (router collapse, expert under-utilization), and (iii) bias the routing toward load uniformity rather than task-optimal assignment. DTR instead routes each token \emph{along depth} over the existing MLLM hidden states, adds no expert networks, no auxiliary loss, no extra communication, and is directly interpretable as token-wise layer preference.

\noindent[\textit{\Rtwo{F1M7}}.] \textbf{\Qtext{Q8: Self-alignment analysis.}} \textbf{\Qtext{A:}} (1) We do not claim removing all pretrained backbones; we remove the dependency on an \emph{external} alignment teacher (DINOv2/CLIP). Unlike REPA-style methods whose teacher is trained purely for representation, our tokenizer is jointly optimized on understanding and generation data, so its target space is intrinsically aligned with the generator and avoids the understanding-generation manifold mismatch external teachers introduce. (2) Self-alignment also accelerates convergence: under our setting, generation loss with self-alignment reaches the no-alignment baseline's final loss using $\sim$20\% fewer training iterations, confirming that internal alignment provides a useful inductive bias even without an external teacher. We will revise the wording.

\noindent[\textit{\Rone{XtYZ}}.] \textbf{\Qtext{Q9: Presentation and citations.}} \textbf{\Qtext{A:}} The Fig.~1(b) caption and the ``self-lignment'' typo are fixed; citation [3] is a blog because FLUX.2-VAE has no paper/arXiv release, and we will add URL plus access date.
\end{document}
